\section{Five rows and larger grids}\label{sec:five}

On five rows the row bound of Theorem~\ref{thm:rowbound} reads $f(5,m)\le m+6$ for
$m\ge3$, and Theorem~\ref{thm:equality} says what a blackout of $m+6$ points has to look
like. Every column is occupied, exactly six columns are doubles, their six row pairs are the
edges of a complete bipartite graph between two rows and the other three, and every other
column is a singleton. Such a blackout passes the axis criterion, so the whole question is the
tilted rectangles. The computed values $f(5,5)=10$, $f(5,6)=11$ and $f(5,7)=12$ of
Section~\ref{sec:verify} fall one short of the bound. This section shows that the bound is
attained at every width from $21$ on. We first prove a finite test for padding any board with
singleton columns in a boundary row, then apply it to one witness of width $21$.

\subsection{A finite test for boundary padding}

Let $h\ge3$ and $m\ge1$, let $S$ be a blackout of the $h\times m$ grid, and call row $1$ or
row $h$ the \emph{padding row}. For integers $L,R\ge0$ we write $S_{L,R}$ for the blackout of
the $h\times(L+m+R)$ grid obtained by adding $L$ columns on the left and $R$ columns on the
right, each new column black exactly in the padding row. We keep the original column indices,
so the added columns are the columns $1-L,\dots,0$ and $m+1,\dots,m+R$, and the black points of
$S_{L,R}$ outside the padding row are exactly the black points of $S$ outside the padding row.
Recall from Section~\ref{sec:graph} that a rectangle is \emph{weak} if it shows at most two
corners, and that $W_h=1+\lfloor(h-1)^2/4\rfloor$ bounds the column span of a tilted rectangle
(Lemma~\ref{lem:span}).

\begin{lemma}[weak tilted rectangles stay near the board]\label{lem:five-local}
Let $h\ge3$.
\begin{enumerate}[(a)]
\item A tilted rectangle has at most one corner in row $1$ and at most one corner in row $h$.
\item In $S_{L,R}$, every weak tilted rectangle has a black corner outside the padding row.
That corner lies in an original column $e$ with $1\le e\le m$, and every column of the
rectangle lies in $[e-W_h,e+W_h]$.
\end{enumerate}
\end{lemma}

\begin{proof}
(a) By Lemma~\ref{lem:leftmost}, from the leftmost corner $(x,j)$ of a tilted rectangle the
two sides are $(a,-b)$ and $(c,d)$ with $a,b,c,d\ge1$. So the four corners lie in the rows
$j$, $j-b$, $j+d$ and $j+d-b$. The smallest of these is $j-b$ and the largest is $j+d$, each
attained once.

(b) A weak tilted rectangle has at least two black corners, and by (a) at most one of them is
in the padding row. So some black corner lies outside the padding row, hence in $S$, hence in a
column $e$ with $1\le e\le m$. By Lemma~\ref{lem:span} the columns of the rectangle differ by
at most $W_h$, so all of them lie in $[e-W_h,e+W_h]$.
\end{proof}

\begin{theorem}[finite collar]\label{thm:collar}
Let $h\ge3$ and $m\ge1$, let $S$ be a blackout of the $h\times m$ grid, let the padding row
be row $1$ or row $h$, and choose the left side, the right side, or both sides. Put
$L^*=W_h$ if the left side is chosen and $L^*=0$ otherwise, and define $R^*$ in the same
way. Then $S_{L,R}$ is valid for every pair $(L,R)$ with $L,R\ge0$, $L=0$ if the left side
is not chosen and $R=0$ if the right side is not chosen, if and only if $S_{L^*,R^*}$ is valid.
\end{theorem}

So one finite board decides every padding length on the chosen sides, with the two lengths
independent. The theorem certifies a proposed padding. It does not say that a given board
passes the test, and Remark~\ref{rem:rightpad} shows one that does not.

\begin{proof}
Write $S^*=S_{L^*,R^*}$ for the test board. If every $S_{L,R}$ is valid, then $S^*$ is
valid, since $(L^*,R^*)$ is one of the admissible pairs. For the converse assume that $S^*$
is valid, fix an admissible pair $(L,R)$, and put $U=S_{L,R}$. We use two facts. First, a
rectangle whose columns lie in the column ranges of both $U$ and $S^*$ has the same corners
and the same presentation in both, because the two blackouts agree on every common column.
Second, by Lemma~\ref{lem:five-local}(b) every weak tilted rectangle of $U$ has its columns
in $[1-W_h,m+W_h]\cap[1-L,m+R]$, and this set is contained in $[1-L^*,m+R^*]$, since
$L=0$ when the left side is not chosen and $R=0$ when the right side is not chosen. So every
weak tilted rectangle of $U$ is a rectangle of $S^*$ with the same presentation.

Suppose that two distinct rectangles $R_1$ and $R_2$ of $U$ have the same presentation $P$.
We derive a contradiction in three cases.

\emph{Step 1. Both axis-aligned.} The graph $G_U$ consists of $G_S$ together with one new
column vertex for each added column, joined by a single edge to the padding row. A vertex of
degree one lies on no cycle, so every cycle of $G_U$ is a cycle of $G_S$, and $G_S$ is a
subgraph of $G_{S^*}$. The board $S^*$ is valid and has $m+W_h\ge3$ columns, so by
Lemma~\ref{lem:girth} the graph $G_{S^*}$ has no cycle of length four or six. Hence $G_U$
has none either, and the other direction of Lemma~\ref{lem:girth}, which needs no hypothesis
on the number of columns, says that no two axis-aligned rectangles of $U$ share a
presentation.

\emph{Step 2. Both tilted.} By Lemma~\ref{lem:threecorners}, $|P|\le2$, so both rectangles
are weak. By the second fact both are rectangles of $S^*$ with the presentation $P$ there,
which contradicts the validity of $S^*$.

\emph{Step 3. One of each.} Let $R_1$ be tilted and $R_2$ axis-aligned. Again $|P|\le2$,
so $R_1$ is weak and lies in the column range of $S^*$. Consider either column of $R_2$. If
it contains a point of $P$, it is a column of $R_1$. Otherwise both corners of $R_2$ in that
column are black, so the column contains two black points of $U$, and it is an original
column, because an added column contains exactly one black point. In both cases the column
lies in the column range of $S^*$. So $R_2$ is also a rectangle of $S^*$ with presentation
$P$, which again contradicts the validity of $S^*$.

Every case contradicts the assumption, so $U$ is valid.
\end{proof}

\begin{remark}[the padding row has to be a boundary row]\label{rem:interior}
Step 3 uses only that an added column holds a single black point, but Step 2 uses
Lemma~\ref{lem:five-local}(a), and the theorem fails for an interior padding row. Take
$h=3$, the original columns $1$, $2$ and $3$ with black rows $\{1\}$, $\{1\}$ and
$\{1,3\}$, and the padding row $2$. Here $W_3=2$. With the two added columns $0$ and $-1$
the board is valid. With a third added column $-2$ the diamond
$\{(-2,2),(-1,1),(-1,3),(0,2)\}$ has both corners of row $2$ black and shows the vertical
pair $\{(-1,1),(-1,3)\}$, and so does the axis-aligned rectangle on the columns $-1$ and $3$
and the rows $1$ and $3$, whose corners in column $3$ are black. So the test board with
$W_3$ added columns passes while a longer padding fails.
\end{remark}

\subsection{The five-row family}

Figure~\ref{fig:witness21} shows the witness of width $21$. Its six double columns are the
columns $1$, $5$, $9$, $13$, $17$ and $21$, black exactly in the row pairs $\{1,2\}$,
$\{1,4\}$, $\{2,3\}$, $\{3,4\}$, $\{2,5\}$ and $\{4,5\}$, and every other column is black
in row $1$ only. The six pairs are the six edges of the complete bipartite graph between the
rows $\{1,3,5\}$ and the rows $\{2,4\}$, so the blackout has the shape that
Theorem~\ref{thm:equality} requires, and it has $21+6=27$ points.

\begin{figure}[t]
\centering
\input{figs/witness21.tex}
\caption{The five-row witness of width $21$. (a) The $27$ black points. The six shaded
columns are the double columns, and every other column is black in row $1$ only. (b) The six
row pairs of the double columns, drawn as a graph on the rows $\{1,3,5\}$ and $\{2,4\}$, each
edge labelled by its column.}
\label{fig:witness21}
\end{figure}

We pad this witness to the left by singleton columns in row $1$. To handle every padding
length at once we work on the half-strip of all columns $x\le21$, with the same six double
columns and the point of row $1$ in every other column, and we call this blackout $S^{\infty}$.
For an integer $c\le1$ let $S_c$ be the restriction of $S^{\infty}$ to the columns
$c,c+1,\dots,21$, a blackout of the $5\times(22-c)$ grid with $22-c+6$ points.

\begin{proposition}[the half-strip is valid]\label{prop:fivecases}
For every integer $c\le1$ the blackout $S_c$ is valid. More precisely, exactly $59$ tilted
rectangles with all corners in columns at most $21$ are weak for $S^{\infty}$. All of them
have their columns in $[-3,21]$, $47$ of them show two corners, $12$ show one corner, none is
hidden, their $59$ presentations are pairwise distinct, and none of them is the presentation
of an axis-aligned rectangle with corners in columns at most $21$.
\end{proposition}

The proof is the hand classification of Appendix~\ref{app:fivecases}. It sorts the tilted
rectangles of a five-row grid into $18$ templates, shows that only the $12$ templates touching
row $1$ can be weak, lists the weak translates of each of them, and checks the presentations
against each other and against the axis-aligned rectangles. Theorem~\ref{thm:collar} gives a
second route. Since $W_5=5$, the board $S_{-4}$ with $26$ columns is the test board for left
padding, and its validity alone implies the validity of every $S_c$. The $3{,}658$ rectangles
of $S_{-4}$ and the $2{,}418$ rectangles of the width-$21$ witness were also enumerated by the
scripts of Section~\ref{sec:verify}, and no two of them share a presentation. These
enumerations are cross-checks of the hand proof, and the proof does not depend on them.

\begin{theorem}[five rows]\label{thm:fiverow}
For every $m\ge21$,
\[
f(5,m)=m+6.
\]
The bound is attained by the blackout whose columns $m-20$, $m-16$, $m-12$, $m-8$, $m-4$ and
$m$ are black exactly in the row pairs $\{1,2\}$, $\{1,4\}$, $\{2,3\}$, $\{3,4\}$, $\{2,5\}$
and $\{4,5\}$, and whose other columns are black in row $1$ only.
\end{theorem}

\begin{proof}
The upper bound is Theorem~\ref{thm:rowbound} with $h=5$ and $m\ge3$. For the lower bound take
$c=22-m\le1$. The blackout $S_c$ of Proposition~\ref{prop:fivecases} occupies the columns
$c,\dots,21$, and moving it $m-21$ columns to the right gives the blackout of the statement
on the columns $1,\dots,m$. Validity is unchanged by the move. The blackout has $m$ occupied
columns and six doubles, so it has $m+6$ points.
\end{proof}

So $f(5,21)=27$, one more than $5+21$, and the excess $f(5,m)-m$ equals $6$ at every width from
$21$ on, which is the largest value the row bound allows. The threshold $21$ is sufficient
and is not claimed to be the smallest, see Proposition~\ref{prop:fivesmall}.

\begin{remark}[padding on the right fails]\label{rem:rightpad}
The family grows to the left, and the direction matters. Append $k$ columns black in row $1$
only to the right of the width-$21$ witness, as columns $22,\dots,21+k$. For $k=1$, $2$ and $3$
the board is valid. For $k=4$ the axis-aligned rectangle $\{(20,1),(20,3),(25,1),(25,3)\}$
and the tilted rectangle $\{(20,3),(21,5),(24,1),(25,3)\}$ both show exactly
$\{(20,3),(25,3)\}$. This agrees with Theorem~\ref{thm:collar}. The test board for the right
side has $W_5=5$ added columns, it contains the board with $k=4$, so it is invalid, and the
theorem then says that some padding length fails. The half turn $(i,j)\mapsto(22-i,6-j)$ of
the width-$21$ board maps the six doubles to themselves and the singletons of row $1$ to
row $5$. So the family turned by a half turn has its doubles at the left end and grows to the
right by singletons in row $5$, and padding it to the left by singletons in row $5$ fails in the
same way.
\end{remark}

\begin{remark}[spacing the doubles further apart does not help]\label{rem:spacing}
Large gaps between the double columns do not make the tilted rectangles safe. Put the same six
row pairs, in the same order, at the columns $1$, $12$, $23$, $34$, $45$ and $56$, with row
$1$ black in every other column of the $5\times56$ grid. Then the axis-aligned rectangle
$\{(41,1),(46,1),(41,3),(46,3)\}$ and the tilted rectangle
$\{(41,3),(42,1),(45,5),(46,3)\}$ both show exactly $\{(41,3),(46,3)\}$. The tilted rectangle
has the horizontal diagonal from $(41,3)$ to $(46,3)$ and the diagonal from $(42,1)$ to
$(45,5)$, of the same length $5$ and the same midpoint. A horizontal diagonal does not force
its partner to be vertical, and the gap $4$ of the witness avoids this coincidence while the
gap $11$ does not.
\end{remark}

\subsection{The \texorpdfstring{$5\times5$}{5 x 5} grid}

The $5\times5$ grid is the first square on which $h+m-1$ is not the answer. The row bound
gives $5+\lfloor25/4\rfloor=11$, and Corollary~\ref{cor:small-m} lowers this to $10$, since
$5<q_5=6$. We give a second proof of the upper bound by counting edges in $G_S$, because it
also describes every optimum.

\begin{lemma}[few edges without short cycles]\label{lem:tenvertices}
Let $G$ be a finite simple bipartite graph with $v$ vertices, $1\le v\le10$, and with no cycle
of length four or six. Then $G$ has at most $v$ edges. If it has exactly $v$ edges, then $G$ is
connected and has exactly one cycle, of length $8$ or $10$.
\end{lemma}

\begin{proof}
Every cycle of a bipartite graph has even length, so every cycle of $G$ has length at least
$8$.

\emph{Step 1. No component has more edges than vertices.} Suppose that a component $G'$ has
more edges than vertices. Remove from $G'$ a vertex of degree one together with its edge, and
repeat until no vertex of degree one is left. Each removal keeps the graph connected and lowers
the numbers of vertices and of edges by one each. So what remains, $H$, is connected and still
has more edges than vertices. In particular $H$ has at least two vertices, so it has no vertex
of degree zero, and every vertex of $H$ has degree at least two. If every degree were exactly
two, $H$ would have as many edges as vertices. So some vertex $x$ has degree at least three.

Let $D_k$ be the set of vertices of $H$ at distance $k$ from $x$. Two different paths from
$x$ to the same vertex, each of length at most $3$, would together contain a cycle of length
at most $6$, which is impossible. So for $k=1,2,3$ every vertex of $D_k$ has exactly one
neighbour in $D_{k-1}$, its parent. It has no neighbour in $D_k$, because $D_k$ lies on one
side of the bipartition. Its degree is at least two, so it has a neighbour in $D_{k+1}$. Two
different vertices of $D_k$ with $k\le2$ cannot share such a neighbour, since that neighbour
would have two parents. So $|D_1|\ge3$, $|D_2|\ge|D_1|\ge3$, $|D_3|\ge|D_2|\ge3$ and
$|D_4|\ge1$. That is at least $1+3+3+3+1=11$ vertices, but $v\le10$. This contradiction
proves Step 1.

\emph{Step 2. Equality.} By Step 1 every component has at most as many edges as vertices, so
$G$ has at most $v$ edges. If $G$ has exactly $v$ edges, every component has exactly as many
edges as vertices, and a connected graph with that property has exactly one cycle. Two
components would mean two disjoint cycles of length at least $8$, which need $16$ vertices.
So $G$ is connected, and its one cycle has even length between $8$ and $10$.
\end{proof}

\begin{theorem}[the $5\times5$ grid]\label{thm:fivefive}
\begin{enumerate}[(a)]
\item $f(5,5)=10$.
\item For every maximum blackout $S$ of the $5\times5$ grid, the graph $G_S$ is connected and
has exactly one cycle, of length $8$ or $10$.
\item (Computer-assisted.) There are exactly $a_5(5)=8$ maximum blackouts. They are the eight
images of the blackout of Figure~\ref{fig:fivefive}(a) under the symmetries of the square,
no two of them coincide, and each has $G_S$ equal to an $8$-cycle with two pendant edges.
\end{enumerate}
\end{theorem}

\begin{proof}
\emph{Upper bound and (b).} If $S$ is valid, then $G_S$ has no cycle of length four or six by
Lemma~\ref{lem:girth}, applied with $m=5\ge3$. The graph is bipartite with $10$ vertices and
$|S|$ edges, so Lemma~\ref{lem:tenvertices} gives $|S|\le10$, and if $|S|=10$ it gives (b).

\emph{Lower bound.} Let $S_0$ be the blackout of Figure~\ref{fig:fivefive}(a), the ten points
\[
(3,1),\ (4,1),\ (5,1),\ (3,2),\ (2,3),\ (3,3),\ (1,4),\ (2,4),\ (1,5),\ (5,5).
\]
Its graph is the $8$-cycle $r_1c_3r_3c_2r_4c_1r_5c_5r_1$ with the two pendant edges $r_2c_3$
and $r_1c_4$, so by Lemma~\ref{lem:girth} no two of the $100$ axis-aligned rectangles share a
presentation. The grid has $130$ rectangles, so $30$ are tilted, and we list them by
Lemma~\ref{lem:leftmost}. From the leftmost corner the sides are $(a,-b)$ and $(c,d)$ with
$ac=bd$, $a+c\le4$ and $b+d\le4$. The case $b+d=2$ gives the nine unit diamonds $D(i,1)$,
$D(i,2)$ and $D(i,3)$ with $1\le i\le3$. The case $b+d=3$ gives the four kinds $A$, $B$, $C$
and $E$ of Lemma~\ref{lem:tilted-narrow}, each with two horizontal and two vertical placements,
sixteen in all. The case $b=d=2$ forces $a=c=2$ and gives the one large diamond
$\{(1,3),(3,1),(5,3),(3,5)\}$. The cases $(b,d)=(1,3)$ and $(3,1)$ give $(a,c)=(1,3)$ or
$(3,1)$, four rectangles that each fill the whole grid. Exactly fifteen of the thirty show at
most two corners, and each of them shows exactly two. They are listed below with the two
visible corners and, in the last column, a corner of the only axis-aligned rectangle that has
both visible points as corners.
\begin{center}\small
\begin{tabular}{@{}rlll@{}}
\toprule
 & tilted rectangle & visible corners & white corner of the rival \\
\midrule
1 & $(1,2),(2,1),(2,3),(3,2)$ & $(1,2),(2,1)$ & $(1,1)$ \\
2 & $(2,2),(3,1),(3,3),(4,2)$ & $(2,2),(4,2)$ & see the text \\
3 & $(3,2),(4,1),(4,3),(5,2)$ & $(4,3),(5,2)$ & $(4,2)$ \\
4 & $(1,3),(2,4),(3,1),(4,2)$ & $(1,3),(4,2)$ & $(1,2)$ \\
5 & $(1,3),(2,2),(2,4),(3,3)$ & $(1,3),(2,2)$ & $(1,2)$ \\
6 & $(1,2),(2,4),(3,1),(4,3)$ & $(1,2),(4,3)$ & $(1,3)$ \\
7 & $(2,3),(3,4),(4,1),(5,2)$ & $(3,4),(5,2)$ & $(5,4)$ \\
8 & $(2,3),(3,2),(3,4),(4,3)$ & $(3,4),(4,3)$ & $(4,4)$ \\
9 & $(2,3),(3,1),(4,4),(5,2)$ & $(4,4),(5,2)$ & $(4,2)$ \\
10 & $(1,4),(2,5),(4,1),(5,2)$ & $(2,5),(5,2)$ & $(2,2)$ \\
11 & $(1,4),(2,5),(3,2),(4,3)$ & $(2,5),(4,3)$ & $(4,5)$ \\
12 & $(1,4),(2,3),(2,5),(3,4)$ & $(2,5),(3,4)$ & $(3,5)$ \\
13 & $(2,4),(3,3),(3,5),(4,4)$ & $(3,5),(4,4)$ & $(3,4)$ \\
14 & $(2,4),(3,2),(4,5),(5,3)$ & $(4,5),(5,3)$ & $(4,3)$ \\
15 & $(2,3),(3,2),(4,5),(5,4)$ & $(4,5),(5,4)$ & $(4,4)$ \\
\bottomrule
\end{tabular}
\end{center}
The fifteen visible pairs are pairwise distinct, so no two of these rectangles share a
presentation, and the other fifteen tilted rectangles show at least three corners and are
determined by Lemma~\ref{lem:threecorners}. It remains to compare the fifteen with the
axis-aligned rectangles. In every row of the table except row $2$ the two visible points lie
in different rows and different columns, so only one axis-aligned rectangle has both as
corners, and the named third corner is white, so that rectangle shows at least three points.
In row $2$ the visible points $(2,2)$ and $(4,2)$ lie in row $2$, and an axis-aligned
rectangle through both uses the columns $2$ and $4$ and one further row. Column $2$ is black
in the rows $3$ and $4$ only and column $4$ in row $1$ only, so no further row is black in
both columns, and no axis-aligned rectangle shows exactly these two points. Hence $S_0$ is
valid and $f(5,5)\ge10$.

\emph{(c).} A maximum blackout has a graph with $10$ edges and no cycle of length four or
six. The script \texttt{verify\_narrow.py} of Section~\ref{sec:verify} lists all $44{,}640$
sets of $10$ points with this property, adding points one at a time and keeping a point only
if no short cycle appears, and checks each of them against all $130$ rectangles. Exactly $8$
are valid, and they are the $8$ images of $S_0$ under the symmetries of the square, all
distinct. The count $8$ was found earlier and independently by the ILP program and by the
check of every blackout from the definition (Table~\ref{tab:values}).
\end{proof}

\begin{figure}[t]
\centering
\input{figs/fivefive.tex}
\caption{The $5\times5$ optimum. (a) The ten black points and the fifteen kept points, which
form the two shaded blocks of five and ten. The bold path turns only at black points and
alternates between rows and columns, and its eight corners are the $8$-cycle of (b). The two
grey points are black points too, and they are the two pendant edges of (b). (b) The graph
$G_S$, an $8$-cycle with two pendant edges.}
\label{fig:fivefive}
\end{figure}

The optimum of Figure~\ref{fig:fivefive}(a) is a staircase of nine points from the top right
corner to the bottom left corner, the pattern of Lemma~\ref{lem:cross}, plus the corner
$(5,5)$. The staircase alone has a tree as its graph, and the extra corner closes a cycle of
length $8$, the shortest the lemma allows. Lemma~\ref{lem:tenvertices} also allows a single
$10$-cycle, that is, two black points in every row and every column. There are $1{,}440$
such sets and none of them is valid. Their graphs have no short cycle, so every collision
involves a tilted rectangle. The count and the exclusion are computed and are not used in any
proof. A
hand proof of the uniqueness in (c) therefore comes down to the blackouts of ten points whose
graph is a connected graph with one cycle of length $8$ or $10$, and to the tilted rectangles
on them.

\subsection{Small widths and the conjecture recorded in the OEIS}

Between the square and the family there is a gap. By Theorem~\ref{thm:equality} a blackout of
$m+6$ points on five rows has every column occupied, six doubles whose row pairs form a complete
bipartite graph between two rows and three rows, and singletons elsewhere. There are
$\binom52=10$ choices of the two rows, $(m)_6=m(m-1)\cdots(m-5)$ injective assignments of
the six pairs to columns and $5^{m-6}$ choices of the singleton rows, that is
$10\,(m)_6\,5^{m-6}$ candidates, which is $7{,}200$ at $m=6$ and $252{,}000$ at $m=7$. All of
them pass the axis criterion, and whether $f(5,m)=m+6$ holds is the question whether one of
them passes the tilted rectangles too.

\begin{proposition}[small five-row widths, computer-assisted]\label{prop:fivesmall}
\begin{enumerate}[(a)]
\item $f(5,5)=10$, $f(5,6)=11$ and $f(5,7)=12$.
\item $f(5,m)\ge m+5$ for $5\le m\le10$ and for every $m\ge17$.
\item {}[[P2-fiverow]]
\end{enumerate}
\end{proposition}

\begin{proof}
(a) The square is Theorem~\ref{thm:fivefive}. For $m=6$ and $m=7$ the ILP program of
Section~\ref{sec:verify} returns the values $11$ and $12$, and the lower bounds $m+5$ are the
stored blackouts of (b). [[P2-fiverow]]

(b) For $5\le m\le10$ the repository stores a valid blackout of $m+5$ points of each size,
found while building the game and rechecked by \texttt{verify\_narrow.py} with its own list of
rectangles. For $m\ge17$ restrict the half-strip of Proposition~\ref{prop:fivecases} to the
columns $18-m,\dots,17$. This board has $m$ columns and the five doubles at the columns
$1$, $5$, $9$, $13$ and $17$, so it has $m+5$ points, and it is valid as the restriction of a
valid blackout. Moved $m-17$ columns to the right, its doubles sit at the columns $m-16$,
$m-12$, $m-8$, $m-4$ and $m$ with the row pairs $\{1,2\}$, $\{1,4\}$, $\{2,3\}$, $\{3,4\}$ and
$\{2,5\}$.

(c) [[P2-fiverow]]
\end{proof}

The entry \seqnum{A397315} records, with the values through the $6\times6$ grid as evidence,
the conjecture that $f(h,m)=h+m$ whenever $\min(h,m)\ge4$, apart from $f(4,4)=7$ and
$f(4,5)=8$. The four-row case and the square $5\times5$ are true (Theorems~\ref{thm:fourvalue}
and~\ref{thm:fivefive}, Corollary~\ref{cor:fourboundary}), and so are the computed cells
$5\times6$, $5\times7$ and $6\times6$. The conjecture fails on five rows.

\begin{theorem}[the conjecture $f(h,m)=h+m$ fails]\label{thm:refute}
$f(5,21)=27>5+21$, and $f(5,m)=m+6=(5+m)+1$ for every $m\ge21$. So the conjecture recorded in
\seqnum{A397315} is false, and the excess $f(5,m)-m$ over the width is $6$, not $5$, at every
width from $21$ on.
\end{theorem}

\begin{proof}
This is Theorem~\ref{thm:fiverow}.
\end{proof}

Three things remain open here. The smallest width $m$ with $f(5,m)=m+6$ is addressed in
Proposition~\ref{prop:fivesmall}(c). Theorem~\ref{thm:periodic} says that at every fixed
height the excess $f(h,m)-m$ is eventually periodic in $m$ with values between $h-1$ and
$q_h$, and on five rows it is eventually constant and equal to its largest possible value
$6$. Whether the same happens at heights six and more, that is, whether $f(h,m)=m+q_h$ for
all large $m$, is open (Section~\ref{sec:open}). On squares, Theorem~\ref{thm:square} shows
that $f(n,n)-2n$ is unbounded, but its threshold is not explicit, and no grid smaller than
$5\times21$ is known to us on which $f$ exceeds $h+m$ (Proposition~\ref{prop:fivesmall}(c)).
